% ============================================================================
%  Capitolo 23 -- Verm1 sul Cartridge Store                      (Lezione 23)
% ============================================================================
\chapter{Verm1 sul Cartridge Store}
\label{ch:store}

\epigraphline{I veri artisti consegnano.}{Steve Jobs, alla squadra del Macintosh, 1983}

\begin{chapgoals}
Alla fine di questo capitolo avrete
\begin{itemize}
\item il vostro \vermone compilato per due architetture, il kit ESP32-C6 e QEMU, dallo stesso sorgente;
\item un manifesto che dà al gioco un'identità: id, titolo, versione, autori, licenza;
\item un'icona, e un colophon con crediti e comandi;
\item un bundle, un solo file ZIP, controllato e pronto per il Cartridge Store;
\item un'idea chiara di ciò che succede dopo la pubblicazione, e di ciò di cui siete responsabili.
\end{itemize}
\end{chapgoals}
% Generated by tools/textbook/export_listings.js from tools/decks/mastery.js and tools/concepts/registry.json: do not edit.
\begin{chapstatus}{Stato di apprendimento di questo capitolo}
\masterychip{now}{DA COMPRENDERE OGGI}\enspace impacchettare e pubblicare; la cartuccia portabile.\par\smallskip
\masterychip{use}{DA UTILIZZARE OGGI}\enspace gli strumenti di PRG32 (si usa, non si studia, in questo corso).\par\smallskip
\masterychip{later}{DA COMPRENDERE PIÙ AVANTI}\enspace niente.
\end{chapstatus}


Questo è l'ultimo capitolo del progetto, ed è di nuovo un tutorial: alla fine
avrete un bundle, \fname{verm1-1.0.0.zip}, controllato e pronto per essere
sottoposto. Tutto fino al caricamento è stato eseguito sul gioco di riferimento
dai test del materiale del corso; il caricamento stesso è un atto personale
compiuto con un token personale, e per questo libro non è stato eseguito.

% ----------------------------------------------------------------------------
\section{Perché pubblicare?}

Un programma che gira solo sul computer del suo autore è un esercizio; uno che
altri possono installare è un prodotto. Pubblicare impone le domande che un
esercizio permette di saltare: come si chiama, chi lo ha fatto, su che cosa
gira, che cosa possono farne gli altri. Il Cartridge Store
(Figura~\ref{fig:st:store}) è lo scaffale di \prg: giochi e programmi che si
installano sul kit, su QEMU e su \prgqt. Ogni cartuccia che vi si trova è stata
scritta da qualcuno che è partito da dove siete partiti voi, e la vostra può
stare accanto a esse. Lo Store è a \url{https://store.prg32.uniparthenope.it},
e \prgqt ha un proprio browser dello Store.

\begin{figure}[htb]
\centering
\includegraphics[width=0.75\textwidth]{prg32qt_store.png}
\caption{Il Cartridge Store, visto da \prgqt.}
\label{fig:st:store}
\end{figure}

Il percorso dal sorgente allo scaffale ha sei stazioni
(Figura~\ref{fig:st:flow}). Le prime quattro avvengono sul computer
dell'autore: niente lo lascia prima del comando che pubblica.

\begin{figure}[htb]
\centering
\begin{tikzpicture}[node distance=3mm]
  \node[blkacc,text width=17mm,minimum height=15mm] (a) {\texttt{verm1.S}\\[1pt]\scriptsize il vostro sorgente};
  \node[blkamb,text width=17mm,minimum height=15mm,right=of a] (b) {2 cartucce\\[1pt]\scriptsize una per architettura};
  \node[blkvio,text width=17mm,minimum height=15mm,right=of b] (c) {metadati\\[1pt]\scriptsize identità, icona, crediti};
  \node[blkgrn,text width=17mm,minimum height=15mm,right=of c] (d) {bundle\\[1pt]\scriptsize un file ZIP};
  \node[blkred,text width=17mm,minimum height=15mm,right=of d] (e) {revisione\\[1pt]\scriptsize un redattore lo controlla};
  \node[blk,fill=prgsteel,text=white,text width=17mm,minimum height=15mm,right=of e] (f) {Store\\[1pt]\scriptsize chiunque installa};
  \foreach \x/\y in {a/b,b/c,c/d,d/e,e/f} \draw[->,thick] (\x) -- (\y);
\end{tikzpicture}
\caption{Dal sorgente allo scaffale.}
\label{fig:st:flow}
\end{figure}

\subsection{Prima di impacchettare: è finito?}

\begin{table}[htb]
\centering\small
\caption{Il gioco è finito?}\label{tab:st:finished}
\begin{tabularx}{\textwidth}{@{}P{0.3\textwidth}LP{0.2\textwidth}@{}}
\toprule
\textbf{Controllo} & \textbf{Come} & \textbf{Con}\\
\midrule
Parte, si gioca e finisce & giocate tre partite intere, perdete in tre modi diversi & \prgqt\\
Nessun frame è in ritardo & telemetria durante una partita lunga: i frame in ritardo restano a zero & \prgqt, modalità debug\\
La memoria è coerente & mettete in pausa a caso e confrontate l'anello, la mappa e lo schermo & \prgqt, vista della memoria\\
Gira sulla seconda architettura & compilate per QEMU, caricate, giocate & QEMU\\
Gira sull'hardware reale & caricate su un kit, giocate con il joystick e l'altoparlante veri & il kit, se ne avete uno\\
È vostro & un titolo, una regola o un aspetto che il riferimento non ha & voi\\
\bottomrule
\end{tabularx}
\end{table}

Il debugger viene usato un'ultima volta, come strumento di qualità
(Tabella~\ref{tab:st:finished}). Una cartuccia che non è mai stata eseguita su
un bersaglio non dovrebbe essere pubblicata per quel bersaglio.

% ----------------------------------------------------------------------------
\section{Passo 1: un sorgente, due cartucce}

\shfile[caption={Compilare per il kit e per QEMU.},label=lst:sto:build]{code/l23/c03.sh}

Il Listato~\ref{lst:sto:build} è il comando di ogni capitolo con un'opzione in
più: \code{-{}-architecture} nomina il bersaglio del file. I nomi seguono una
regola, \fname{gioco-architettura.prg32}. Il compilatore di cartucce risponde
con il nome, l'indirizzo di caricamento \hex{40800000} e la dimensione del
codice: \vermone finito è di circa 6 kilobyte (6\,216 byte per il gioco di
riferimento).

\txtfile[caption={È portabile?},label=lst:sto:summary]{code/l23/c04.sh}

Il comando \code{summary} stampa ciò che il file dichiara di sé
(Listato~\ref{lst:sto:summary}). Il valore \code{abi-table} di
\code{import\_model} significa che la cartuccia raggiunge il runtime attraverso
una tabella di funzioni, la ABI di \prg, di cui \code{prg32\_gfx\_rect} e le
altre sono elementi. Una cartuccia di questo tipo non dipende da una
particolare compilazione del firmware: è \term{portabile}, e lo Store accetta
cartucce portabili. Una cartuccia non portabile è collegata agli indirizzi di
una sola immagine del firmware, e si rompe al successivo aggiornamento.
Ripetete il controllo sul file per QEMU.

% ----------------------------------------------------------------------------
\section{Passo 2: un'identità}

\subsection{\texttt{manifest.json}: chi è questo gioco?}

Create il file del Listato~\ref{lst:sto:manifest} in
\fname{\textasciitilde/verm1} con un editor di testo. (Nel listato la
descrizione è tagliata; il file completo del gioco di riferimento è
nell'Appendice~\ref{app:verm1}.)

\jsonfile[caption={\fname{manifest.json}.},label=lst:sto:manifest]{code/l23/c05.json}

La riga \code{abi} dice quale versione del formato dei metadati è questa. Gli
altri campi sono impegni (Tabella~\ref{tab:st:fields}).

\begin{table}[htb]
\centering\small
\caption{I campi del manifesto, e a che cosa impegnano il loro autore.}\label{tab:st:fields}
\begin{tabularx}{\textwidth}{@{}lLL@{}}
\toprule
\textbf{Campo} & \textbf{Regola} & \textbf{Per \vermone}\\
\midrule
\code{id} & per consuetudine un nome di dominio rovesciato; non cambia mai & uno per autore: metteteci il vostro nome\\
\code{version} & major.minor.patch; una versione pubblicata non viene mai sostituita & 1.0.0, poi 1.0.1 per una correzione\\
\code{title} & ciò che il giocatore legge & Verm1, o il nome della vostra variante\\
\code{authors} & le persone che lo hanno scritto & voi, e chiunque abbia lavorato con voi\\
\code{license} & che cosa possono farne gli altri & MIT, se volete che altri imparino da esso\\
\code{tags} & come verrà trovato & arcade, snake, assembly, education\\
\bottomrule
\end{tabularx}
\end{table}

\begin{prgwarn}
Trattate un id e una versione, una volta pubblicati, come permanenti. Un errore
si corregge con una nuova versione, non con una sostituzione. E due autori non
possono pubblicare lo stesso id: ciascuno usa un id proprio, per esempio
\code{it.uniparthenope.arc12.}\emph{cognome}\code{.verm1}.
\end{prgwarn}

\subsection{Un'icona}

L'icona è un piccolo PNG che rappresenta il gioco sullo scaffale
(Figura~\ref{fig:st:icon}). L'icona di riferimento è di 64 per 64 pixel: un
verme e un pezzo di cibo, nei colori del gioco. La pixel art si addice a una
macchina di 320 per 200 pixel: disegnatela in grande, cella per cella, con un
qualsiasi editor che salvi PNG, oppure disegnatela con un programma, come è
stata fatta l'icona di riferimento. Salvatela come \fname{icon.png} accanto al
manifesto, che la nomina sotto \code{assets}. Un'icona presa dal web è il
lavoro di qualcun altro: disegnate la vostra.

\begin{figure}[htb]
\centering
\includegraphics[width=28mm]{verm1_icon.png}
\caption{\fname{icon.png}, ingrandita.}
\label{fig:st:icon}
\end{figure}

\subsection{\texttt{colophon.json}: crediti e comandi}

\jsonfile[caption={\fname{colophon.json}.},label=lst:sto:colophon]{code/l23/c06.json}

\emph{Colophon} è la parola dei tipografi per la nota alla fine di un libro che
dice chi lo ha fatto. Qui è la pagina dei crediti della cartuccia
(Listato~\ref{lst:sto:colophon}): ruoli e nomi, dando credito a tutti; i
comandi, cioè che cosa fa ciascun ingresso, per il giocatore; i
ringraziamenti. Sostituite dappertutto \texttt{Your Name}, e impostate la data
di rilascio. Il primo ringraziamento resta: un gioco che parte
dall'implementazione di riferimento di \vermone, contributo di Giuseppe
Coviello, lo dice.

\subsection{Un omaggio non è una copia}

\vermone è un omaggio a Snake: l'idea di un verme che cresce è vecchia quanto
\emph{Blockade} (1976) ed è libera per tutti. Il nome, il codice, l'icona e i
suoni sono originali: niente è stato preso da un altro gioco. Le regole di un
gioco non sono protette; il codice, le immagini, la musica e i nomi lo sono.
Prima di pubblicare, chiedetevi per ogni pezzo: l'ho scritto o disegnato io?
Se no, la sua licenza mi permette di pubblicarlo, e ho dato credito? I nomi e
i personaggi dei giochi commerciali sono marchi, e non vanno usati. La licenza
scelta dice agli altri che cosa possono fare con l'opera: la MIT permette di
studiarla, modificarla e condividerla, conservando il nome dell'autore.

\begin{prgconcept}[La regola]
Sopra c'è il vostro nome. Pubblicate solo ciò che firmereste.
\end{prgconcept}

% ----------------------------------------------------------------------------
\section{Passo 3: la cartuccia porta con sé i suoi documenti}

\shfile[caption={\fname{metadata.json}, e come allegarlo.},label=lst:sto:attach]{code/l23/c07.sh}

Il file \fname{metadata.json} è una copia del manifesto senza le due voci
\code{assets} e \code{architectures}. Il comando \code{attach-metadata} scrive
titolo, icona e crediti \emph{dentro} la cartuccia, una volta per file e con la
sua architettura (Listato~\ref{lst:sto:attach}): una cartuccia copiata a mano
dice ancora che cosa è. Lo strumento risponde, per il gioco di riferimento,
\code{legacy=6316 blocks=3}: i tre blocchi sono i metadati, l'icona e il
colophon, accodati dopo il programma, che di per sé non viene toccato.

\begin{prgwarn}
Dare il manifesto completo a \code{-{}-metadata} fa fallire lo strumento, a
causa della voce \code{assets}.
\end{prgwarn}

\txtfile[caption={Rileggere i metadati (uscita ridotta).},label=lst:sto:inspect]{code/l23/c08.sh}

\begin{prgconcept}
Non fidatevi mai di una scrittura che non avete riletto. È l'abitudine di tutto
il libro: la vista della memoria del debugger, e ora questo.
\end{prgconcept}

Controllate l'id, il titolo, la versione e l'architettura
(Listato~\ref{lst:sto:inspect}); \code{trailer\_present} dice che i documenti
sono nel file. Poi eseguite ancora una volta la cartuccia: deve funzionare
ancora.

% ----------------------------------------------------------------------------
\section{Passo 4: il bundle}

\txtfile[caption={Impacchettare il bundle (uscita reale per il gioco di riferimento).},label=lst:sto:pack]{code/l23/c09.sh}

Lo strumento legge il manifesto e raccoglie i file che esso nomina
(Listato~\ref{lst:sto:pack}), che devono essere nella stessa cartella del
manifesto. Un file mancante o con il nome sbagliato lo ferma qui, e non allo
Store. Il nome del bundle porta la versione. L'intero gioco, per due macchine,
è di 16 kilobyte.

\begin{prgcheck}[Eseguite: provate il bundle, non la vostra cartella]
\begin{enumerate}[nosep]
\item Decomprimete \fname{verm1-1.0.0.zip} in una cartella vuota.
\item In \prgqt importate \fname{verm1-esp32c6.prg32} da quella cartella ed
  eseguitelo: titolo, icona e crediti devono essere i vostri.
\item Caricate \fname{verm1-qemu.prg32} su QEMU e giocate una partita.
\item Se avete un kit, caricate il file per l'ESP32-C6 e giocate con il
  joystick vero.
\item Chiedete a un collega di fare lo stesso su un altro computer: il primo
  utente che non siete voi.
\end{enumerate}
\textbf{Risultato atteso.} Il gioco gira dai file del bundle, su ogni bersaglio
dichiarato, su un computer che non è il vostro. Provare i file impacchettati,
e non quelli della cartella di lavoro, coglie l'errore classico: un bundle
costruito prima dell'ultima correzione.
\end{prgcheck}

% ----------------------------------------------------------------------------
\section{Passo 5: la pubblicazione}

\shfile[caption={Pubblicare: un token e un comando.},label=lst:sto:publish]{code/l23/c11.sh}

Per pubblicare, l'autore accede allo Store, crea un \term{token} nella pagina
\emph{Tokens}, lo tiene nell'ambiente della shell e sottopone il bundle
(Listato~\ref{lst:sto:publish}).

\begin{prgwarn}[Il token siete voi]
Chi ha il token pubblica a vostro nome. È un segreto: non va nel sorgente, in
un file, in un repository o in una schermata. Se viene perso o mostrato per
errore, revocatelo e createne uno nuovo.
\end{prgwarn}

\subsection{Che cosa succede dopo}

\begin{description}[style=nextline,leftmargin=1.5em]
\item[In attesa.] Il bundle è memorizzato come proposta: non è ancora sullo
  scaffale.
\item[Revisione.] Un redattore controlla che si installi, che giri, che sia
  descritto onestamente e che non contenga niente che non sia dell'autore.
\item[Pubblicato.] Il gioco compare nello Store, per le architetture
  dichiarate.
\item[Un errore?] Correggetelo, portate la versione a 1.0.1, impacchettate e
  sottoponete di nuovo: la 1.0.0 resta com'era.
\end{description}

La moderazione protegge i giocatori e gli autori. Un rifiuto arriva con una
motivazione, e gli si risponde con una nuova versione.

\subsection{La lista di controllo per la sottomissione}

Una lista di controllo non è burocrazia: è il modo in cui piloti e chirurghi
evitano gli errori che fanno gli esperti.

\begin{enumerate}
\item Il gioco è stato giocato fino alla fine su ogni architettura del
  manifesto.
\item \code{cartridge summary} dice \code{import\_model: abi-table} per ogni
  file.
\item L'id è vostro e unico; la versione è nuova; titolo, sommario e
  descrizione dicono che cosa è il gioco.
\item Autori, licenza e colophon sono completi; ogni pezzo è vostro o
  accreditato.
\item \code{inspect-metadata} mostra l'id, la versione e l'architettura giusti
  in ogni cartuccia.
\item Il bundle è stato decompresso e provato su un altro computer.
\item Il token non è in alcun file, in alcuna schermata e in alcun repository.
\end{enumerate}

% ----------------------------------------------------------------------------
% Generated by tools/textbook/export_listings.js from tools/concepts/assessments.json: do not edit.
\begin{chapstatus}{Il filo del capitolo}
Un gioco è finito quando qualcun altro può eseguirlo: una cartuccia portabile, metadati onesti, un bundle provato su un altro computer e un token tenuto segreto. Prossimo passo: macchine più veloci di un solo core.
\par\medskip\textbf{Biglietto d'uscita} (da due a quattro minuti, da soli, su carta):
\begin{enumerate}[nosep,leftmargin=2em]
\item cartridge summary stampa import\_model: abi-table. Che cosa dice questo della cartuccia, e perché lo Store accetta solo cartucce del genere?
\item Avete pubblicato la versione 1.0.0 e trovato un errore. Che cosa pubblicate ora? Dove non deve mai comparire il token dello Store?
\end{enumerate}
\end{chapstatus}


\section{Verificate di aver capito}

\begin{quiz}{Avete pubblicato la versione 1.0.0 e trovato un errore. Che cosa fate?}
\item caricate di nuovo la 1.0.0
\item pubblicate la 1.0.1
\item cambiate l'id
\item chiedete che il vecchio file venga modificato
\end{quiz}

\begin{quiz}{Perché \fname{metadata.json} non contiene la voce \code{architectures}?}
\item è facoltativa
\item i metadati vanno dentro una cartuccia, che ha una sola architettura
\item QEMU non la supporta
\item sarebbe troppo lunga
\end{quiz}

\begin{quiz}{Dove va tenuto il token dello Store?}
\item in \fname{verm1.S}, come commento
\item nel manifest
\item in una variabile d'ambiente della vostra shell
\item nel \fname{README} del repository
\end{quiz}

\section{Il vostro rilascio}

\begin{exercise}[l'elaborato del progetto]
\begin{enumerate}[nosep]
\item Scegliete il vostro id e il vostro titolo, e scrivete il manifesto e il
  colophon per la vostra variante di \vermone.
\item Disegnate la vostra icona.
\item Compilate, allegate, ispezionate, impacchettate: producete il vostro
  bundle, e superate la lista di controllo con un collega come revisore.
\item Se volete il vostro gioco sullo scaffale, pubblicatelo.
\end{enumerate}
\expected{Un bundle del vostro gioco, rivisto da un collega rispetto alla
lista di controllo. La pubblicazione è facoltativa e personale; il bundle, con
il sorgente della vostra variante, è l'elaborato del progetto, e l'esame orale
può partire da una qualsiasi delle sue righe.}
\end{exercise}

\section{Che cosa avete costruito}

Un gioco di seicento righe di assembly, a partire da tre funzioni che
ritornavano soltanto. Con esso, ogni strato della macchina: bit e numeri, porte
e tabelle, stato e memoria, il processore, lo schermo, i dispositivi. Un
metodo: un passo eseguibile alla volta, e uno sguardo dentro ogni volta che
qualcosa sorprende. E l'ultimo passo, che ne fa ingegneria: un prodotto con un
nome, una versione e un autore.

Il prossimo gioco sta a voi sceglierlo. Il prossimo capitolo mostra dove le
stesse istruzioni girano in grande scala.
